% ============================================================
%  GUÍA 4: Funciones Racionales, Composición e Inversas
%  Curso de Introducción a la Universidad — Precálculo y Cálculo I
%  Autor: Ing. Gabriel Astudillo
%  Clase cubierta: 4
% ============================================================

\documentclass[12pt, a4paper]{article}

% ── Paquetes ─────────────────────────────────────────────────
\usepackage{mathwizards-guia}
\mwguia{Guía 4 — Funciones III}{Ing. Gabriel Astudillo $\cdot$ Curso Preuniversitario}

\begin{document}

% ── Portada / Título ─────────────────────────────────────────
\mwportada{Guía de Estudio 4}{Funciones Racionales, Composición e Inversas}{Asíntotas, Discontinuidades, Álgebra de Funciones, Composición, Inyectividad e Inversas}

\vspace{4pt}
\begin{cuadroinfo}
  \small
  \textbf{Presentación:} Esta guía profundiza en las funciones racionales —asíntotas verticales, horizontales y oblicuas, discontinuidades evitables— y en el álgebra de funciones: suma, producto, cociente, composición y función inversa, con el criterio de inyectividad y la prueba de la recta horizontal. La teoría, los ejemplos resueltos y los ejercicios propuestos con clave te preparan para los límites y la derivación.
  \\[4pt]
  \textbf{Nivel de Dificultad:}\quad
  \facil{} Básico / Directo \quad
  \medio{} Intermedio \quad
  \dificil{} Desafiante \quad
  \muyDificil{} Muy Difícil / Reto
\end{cuadroinfo}

\vspace{8pt}

\section{Funciones Racionales y Asíntotas}
% ============================================================

\subsection{Conceptos y técnicas}

\begin{cuadroteoria}
  \textbf{Función racional:} $f(x) = \dfrac{P(x)}{Q(x)}$, donde $P$ y $Q$ son polinomios y $Q(x) \neq 0$.

  \textbf{Dominio:} Todos los $x$ tales que $Q(x) \neq 0$.

  \textbf{Asíntotas verticales:} En $x = a$ si $Q(a) = 0$ y $P(a) \neq 0$. Si ambos se anulan, hay un \textbf{hueco} (discontinuidad evitable) en $x = a$.

  \textbf{Asíntotas horizontales:}
  \begin{itemize}[leftmargin=*]
    \item Si grado$(P) <$ grado$(Q)$: $y = 0$.
    \item Si grado$(P) =$ grado$(Q)$: $y = \dfrac{\text{coef. principal de } P}{\text{coef. principal de } Q}$.
    \item Si grado$(P) >$ grado$(Q)$: no hay asíntota horizontal (puede haber oblicua si la diferencia es 1).
  \end{itemize}

  \textbf{Asíntota oblicua:} Si grado$(P) =$ grado$(Q) + 1$, se obtiene al dividir $P$ entre $Q$; el cociente es la asíntota.
\end{cuadroteoria}

\newpage

\subsection{Ejercicios resueltos}

\textbf{Ejemplo R1.} Encuentra dominio, asíntotas y huecos de $f(x) = \dfrac{x^2 - 1}{x^2 - 3x + 2}$.

\begin{cuadrosolucion}
  \textbf{Solución:}
  \begin{enumerate}
    \item Factorizamos: $f(x) = \dfrac{(x - 1)(x + 1)}{(x - 1)(x - 2)} = \dfrac{x + 1}{x - 2}$ (para $x \neq 1$).
    \item \textbf{Dominio:} $(-\infty, 1) \cup (1, 2) \cup (2, \infty)$.
    \item \textbf{Hueco:} En $x = 1$, porque el factor $(x - 1)$ se cancela. El valor de $y$ en el hueco: $f(1) = \frac{2}{-1} = -2$. Hueco en $(1, -2)$.
    \item \textbf{Asíntota vertical:} En $x = 2$ (no se cancela con el numerador).
    \item \textbf{Asíntota horizontal:} Ambos grados son 2. Cociente de coeficientes principales: $y = \frac{1}{1} = 1$.
  \end{enumerate}
\end{cuadrosolucion}

\textbf{Ejemplo R2.} Encuentra las asíntotas de $f(x) = \dfrac{2x^2 + 3x - 1}{x - 2}$.

\begin{cuadrosolucion}
  \textbf{Solución:} Grado del numerador (2) es uno más que el del denominador (1), entonces hay asíntota oblicua. Dividimos:
  $$\dfrac{2x^2 + 3x - 1}{x - 2} = 2x + 7 + \dfrac{13}{x - 2}.$$
  Asíntota oblicua: $y = 2x + 7$. Asíntota vertical: $x = 2$. No hay horizontal.
\end{cuadrosolucion}

\textbf{Ejemplo R3.} Encuentra dominio y asíntotas de $f(x) = \dfrac{3}{x^2 + 1}$.

\begin{cuadrosolucion}
  \textbf{Solución:} $x^2 + 1$ nunca es cero, así que el dominio es $\mathbb{R}$. No hay asíntotas verticales. Grado numerador (0) $<$ grado denominador (2), así que $y = 0$ es asíntota horizontal.
\end{cuadrosolucion}

\newpage

\textbf{Ejemplo R4.} Grafica $f(x) = \dfrac{2x + 1}{x - 3}$ indicando asíntotas e interceptos.

\begin{cuadrosolucion}
  \textbf{Solución:}
  \begin{itemize}
    \item Asíntota vertical: $x = 3$.
    \item Asíntota horizontal: ambos grados 1, coeficientes $2/1$: $y = 2$.
    \item Intercepto $y$: $f(0) = \frac{1}{-3} = -\frac{1}{3}$; punto $(0, -\frac{1}{3})$.
    \item Intercepto $x$: $2x + 1 = 0 \Rightarrow x = -\frac{1}{2}$; punto $(-\frac{1}{2}, 0)$.
  \end{itemize}
  \begin{center}
    \begin{tikzpicture}[scale=0.7]
      \clip (-6.6,-5.6) rectangle (7.6,6.6);
      \draw[thin, gray!40] (-6,-5) grid (7,6);
      \draw[thick, -Latex] (-6.5,0) -- (7.5,0) node[right] {$x$};
      \draw[thick, -Latex] (0,-5.5) -- (0,6.5) node[above] {$y$};
      \draw[dashed, red!80, thick] (3,-5.5) -- (3,6.5) node[right] {$x=3$};
      \draw[dashed, red!80, thick] (-6,2) -- (7,2) node[right] {$y=2$};
      \draw[smooth, very thick, mwprimario, domain=-6:2.4, samples=100] plot (\x, {(2*\x+1)/(\x-3)});
      \draw[smooth, very thick, mwprimario, domain=3.6:7, samples=100] plot (\x, {(2*\x+1)/(\x-3)});
      \filldraw[black] (0,-1/3) circle (2.5pt);
      \filldraw[black] (-0.5,0) circle (2.5pt);
    \end{tikzpicture}
  \end{center}
\end{cuadrosolucion}

\subsection{Ejercicios propuestos}

\begin{enumerate}[series=ejercicios, leftmargin=*]
  \item \facil{} Encuentra el dominio de $f(x) = \dfrac{x + 2}{x - 5}$.

  \item \facil{} Encuentra el dominio de $f(x) = \dfrac{x^2}{x^2 - 9}$.

  \item \facil{} Encuentra las asíntotas de $f(x) = \dfrac{1}{x - 2}$.

  \item \facil{} Encuentra las asíntotas de $f(x) = \dfrac{5x + 1}{x}$.

  \item \medio{} Encuentra dominio, huecos, asíntotas e interceptos de $f(x) = \dfrac{x^2 - 4}{x^2 - 2x - 3}$.

  \item \medio{} Encuentra las asíntotas de $f(x) = \dfrac{2x^2 + 5x + 3}{x - 1}$.

  \item \medio{} Encuentra las asíntotas de $f(x) = \dfrac{x^2}{x^2 + 1}$.

  \item \medio{} Determina si $f(x) = \dfrac{x^3 - 8}{x - 2}$ tiene hueco o asíntota en $x = 2$ y simplifica la función.

  \item \dificil{} Encuentra la asíntota oblicua de $f(x) = \dfrac{3x^3 + 2x^2 - x + 1}{x^2 + 1}$.

  \item \dificil{} Determina los valores de $k$ para que $f(x) = \dfrac{kx + 1}{x - 2}$ tenga asíntota horizontal en $y = 4$.

  \item \dificil{} Una función racional tiene asíntota vertical en $x = -2$, asíntota horizontal en $y = 3$ e intercepto $y$ en $(0, \frac{3}{2})$. Encuentra una posible fórmula para $f(x)$.

  \item \dificil{} Grafica $f(x) = \dfrac{x^2 - 1}{(x - 2)^2}$ indicando todas sus asíntotas, huecos e interceptos.
\end{enumerate}


% ============================================================

\section{Composición de Funciones y Función Inversa}
% ============================================================

\subsection{Operaciones y composición}

\begin{cuadroteoria}
  Dadas $f$ y $g$:
  \begin{align*}
    (f + g)(x)                  & = f(x) + g(x)                          \\
    (f - g)(x)                  & = f(x) - g(x)                          \\
    (fg)(x)                     & = f(x) \cdot g(x)                      \\
    \left(\frac{f}{g}\right)(x) & = \frac{f(x)}{g(x)}, \quad g(x) \neq 0
  \end{align*}
\end{cuadroteoria}

\begin{cuadrosolucion}
  \textbf{Composición de funciones:}
  $$(f \circ g)(x) = f(g(x))$$
  Es decir, evaluamos primero $g$ en $x$ y luego evaluamos $f$ en el resultado. \textbf{El orden importa:} en general, $f \circ g \neq g \circ f$.

  El dominio de $f \circ g$ es el conjunto de $x$ en el dominio de $g$ tales que $g(x)$ está en el dominio de $f$.
\end{cuadrosolucion}

\newpage

\subsection{Función inversa}

\begin{cuadroteoria}
  \textbf{Definición:} Una función $f$ con dominio $A$ y rango $B$ tiene inversa $f^{-1}$ si se cumple:
  $$f^{-1}(f(x)) = x \quad \text{para todo } x \in A, \qquad f(f^{-1}(y)) = y \quad \text{para todo } y \in B.$$

  \textbf{Existencia:} $f$ tiene inversa si y solo si $f$ es \textbf{inyectiva} (uno a uno), es decir, si $f(x_1) = f(x_2) \Rightarrow x_1 = x_2$. Gráficamente: la \textbf{prueba de la recta horizontal} debe pasar.

  \textbf{Método para hallar la inversa:}
  \begin{enumerate}[leftmargin=*]
    \item Escribe $y = f(x)$.
    \item Intercambia $x$ y $y$.
    \item Despeja $y$: el resultado es $y = f^{-1}(x)$.
  \end{enumerate}

  \textbf{Propiedad gráfica:} Las gráficas de $f$ y $f^{-1}$ son simétricas respecto a la recta $y = x$.
\end{cuadroteoria}

\begin{cuadroadvertencia}
  \textbf{Restricción del dominio:} Si una función no es inyectiva en todo su dominio, se puede \textbf{restringir} el dominio para que lo sea. Por ejemplo, $f(x) = x^2$ no es inyectiva en $\mathbb{R}$, pero sí lo es en $[0, \infty)$ (en ese caso su inversa es $f^{-1}(x) = \sqrt{x}$).
\end{cuadroadvertencia}

\subsection{Ejercicios resueltos}

\textbf{Ejemplo R1.} Dadas $f(x) = 2x + 3$ y $g(x) = x^2$:
\begin{enumerate}[label=\alph*)]
  \item $(f \circ g)(x)$
  \item $(g \circ f)(x)$
  \item $(f \circ g)(-2)$
\end{enumerate}

\begin{cuadrosolucion}
  \textbf{Solución:}
  \begin{itemize}
    \item $(f \circ g)(x) = f(g(x)) = f(x^2) = 2x^2 + 3$.
    \item $(g \circ f)(x) = g(f(x)) = g(2x + 3) = (2x + 3)^2 = 4x^2 + 12x + 9$.
    \item $(f \circ g)(-2) = 2(-2)^2 + 3 = 2(4) + 3 = 11$.
  \end{itemize}
  \textbf{Observa que $f \circ g \neq g \circ f$.}
\end{cuadrosolucion}

\newpage

\textbf{Ejemplo R2.} Dadas $f(x) = \dfrac{1}{x - 3}$ y $g(x) = \sqrt{x}$, calcula $(f \circ g)(x)$ y determina su dominio.

\begin{cuadrosolucion}
  \textbf{Solución:}
  \begin{enumerate}
    \item $(f \circ g)(x) = f(g(x)) = f(\sqrt{x}) = \dfrac{1}{\sqrt{x} - 3}$.
    \item Para el dominio necesitamos que $x$ esté en el dominio de $g$ \textbf{y} que $g(x)$ esté en el dominio de $f$:
          \begin{itemize}
            \item Dominio de $g$: $x \geq 0$ (raíz cuadrada).
            \item $g(x)$ en el dominio de $f$: $\sqrt{x} \neq 3$, es decir, $x \neq 9$.
          \end{itemize}
    \item Dominio de $f \circ g$: $[0, 9) \cup (9, \infty)$.
  \end{enumerate}
\end{cuadrosolucion}

\textbf{Ejemplo R3.} Halla $f^{-1}(x)$ para $f(x) = \dfrac{2x - 1}{x + 3}$.

\begin{cuadrosolucion}
  \textbf{Solución:}
  \begin{enumerate}
    \item $y = \dfrac{2x - 1}{x + 3}$.
    \item Intercambiamos: $x = \dfrac{2y - 1}{y + 3}$.
    \item Despejamos $y$: $x(y + 3) = 2y - 1 \Rightarrow xy + 3x = 2y - 1 \Rightarrow xy - 2y = -3x - 1 \Rightarrow y(x - 2) = -3x - 1 \Rightarrow y = \dfrac{-3x - 1}{x - 2} = \dfrac{3x + 1}{2 - x}$.
  \end{enumerate}
  Por tanto: $f^{-1}(x) = \dfrac{3x + 1}{2 - x}$, con dominio $x \neq 2$.
\end{cuadrosolucion}

\textbf{Ejemplo R4.} Determina si $f(x) = x^3 + 1$ es inyectiva y halla su inversa.

\begin{cuadrosolucion}
  \textbf{Solución:} $f$ es inyectiva porque es estrictamente creciente (o porque $f(x_1) = f(x_2)$ implica $x_1^3 = x_2^3$, luego $x_1 = x_2$). Para hallar la inversa:
  \begin{enumerate}
    \item $y = x^3 + 1$.
    \item Intercambiamos: $x = y^3 + 1 \Rightarrow y^3 = x - 1$.
    \item $y = \sqrt[3]{x - 1}$.
  \end{enumerate}
  $f^{-1}(x) = \sqrt[3]{x - 1}$.
\end{cuadrosolucion}

\newpage

\textbf{Ejemplo R5.} Para $f(x) = x^2 + 4x$ con dominio restringido a $[-2, \infty)$, encuentra $f^{-1}(x)$.

\begin{cuadrosolucion}
  \textbf{Solución:}
  \begin{enumerate}
    \item $y = x^2 + 4x$.
    \item Intercambiamos: $x = y^2 + 4y \Rightarrow y^2 + 4y - x = 0$.
    \item Resolvemos la cuadrática en $y$: $y = \dfrac{-4 \pm \sqrt{16 + 4x}}{2} = -2 \pm \sqrt{x + 4}$.
    \item Elegimos el signo $+$ porque el dominio restringido es $x \geq -2$, donde la rama creciente corresponde a $+$:
          $f^{-1}(x) = -2 + \sqrt{x + 4}$.
  \end{enumerate}
\end{cuadrosolucion}

\subsection{Ejercicios propuestos}

\begin{enumerate}[resume=ejercicios, leftmargin=*]
  \item \facil{} Dadas $f(x) = x + 2$ y $g(x) = 3x$, calcula:
        \begin{enumerate}[label=\alph*)]
          \item $(f + g)(x)$
          \item $(f - g)(x)$
          \item $(fg)(x)$
          \item $\left(\frac{f}{g}\right)(x)$
        \end{enumerate}

  \item \facil{} Dadas $f(x) = x^2$ y $g(x) = x + 1$, calcula $(f \circ g)(x)$ y $(g \circ f)(x)$.

  \item \facil{} Halla la inversa de $f(x) = 3x - 7$.

  \item \facil{} Halla la inversa de $f(x) = x^3 + 5$.

  \item \medio{} Dadas $f(x) = \sqrt{x + 1}$ y $g(x) = x^2$, calcula $(f \circ g)(x)$ y su dominio.

  \item \medio{} Determina si $f(x) = 2x^2 - 8x + 5$ es inyectiva en todo $\mathbb{R}$. Si no lo es, indica una restricción del dominio adecuada para que sea inyectiva.

  \item \medio{} Halla la inversa de $f(x) = \dfrac{x}{x - 2}$.

  \item \medio{} Verifica que $f$ y $g$ son inversas entre sí calculando $f(g(x))$ y $g(f(x))$:
        $$f(x) = \frac{3x - 1}{2}, \qquad g(x) = \frac{2x + 1}{3}.$$

  \item \dificil{} Sea $f(x) = x^2 - 4x + 3$ con dominio restringido a $[2, \infty)$. Encuentra $f^{-1}(x)$.

  \item \dificil{} Dadas $f(x) = \dfrac{1}{x - 1}$ y $g(x) = x^2$, encuentra el dominio de $(f \circ g)(x)$ y de $(g \circ f)(x)$.

  \item \dificil{} Si $f(g(x)) = 2x - 1$ donde $g(x) = x + 3$, encuentra $f(x)$.

  \item \dificil{} Un impuesto de $20\%$ se aplica al precio de un artículo, y luego se aplica un descuento de \$10. Si el precio final es $P(x) = 0{,}8x - 10$ (donde $x$ es el precio original), encuentra $P^{-1}(x)$ e interpreta su significado en el contexto.
\end{enumerate}

\newpage

% ============================================================

% ──────────────────────────────────────────────────────────
%  NUEVOS EJERCICIOS PROPUESTOS (26) — INSERCIÓN MANUAL
% ──────────────────────────────────────────────────────────
\subsection{Ejercicios propuestos: Asíntotas y Gráficas Racionales}

\begin{enumerate}[resume=ejercicios, leftmargin=*]
  \item \facil{} Determina el dominio de $f(x) = \dfrac{x + 2}{x^2 - 9}$.

  \item \facil{} Halla la asíntota vertical de $f(x) = \dfrac{2x + 1}{x - 4}$.

  \item \medio{} Determina las asíntotas vertical y horizontal de $f(x) = \dfrac{3x + 6}{x - 1}$.

  \item \medio{} Halla la asíntota horizontal de $f(x) = \dfrac{2x^2 - 1}{3x^2 + 5x}$.

  \item \medio{} Clasifica la discontinuidad de $f(x) = \dfrac{x^2 - 4}{x - 2}$ y escribe la función simplificada equivalente.

  \item \dificil{} Determina las asíntotas vertical y oblicua de $f(x) = \dfrac{x^2 + 1}{x - 1}$.

  \item \dificil{} Halla la asíntota oblicua de $f(x) = \dfrac{2x^2 - 3x + 1}{x + 2}$.

  \item \dificil{} Describe el bosquejo de $f(x) = \dfrac{x}{x^2 - 4}$: dominio, asíntotas e interceptos.

  \item \muyDificil{} Determina el dominio, las asíntotas verticales y la asíntota oblicua de $f(x) = \dfrac{x^3}{x^2 - 1}$.
\end{enumerate}

\subsection{Ejercicios propuestos: Composición e Inversas Avanzadas}

\begin{enumerate}[resume=ejercicios, leftmargin=*]
  \item \facil{} Dadas $f(x) = 2x + 1$ y $g(x) = x^2$, halla $(f \circ g)(x)$ y $(g \circ f)(x)$.

  \item \facil{} Halla la función inversa de $f(x) = 3x - 6$.

  \item \medio{} Halla la función inversa de $f(x) = \sqrt{x - 2}$ e indica su dominio.

  \item \medio{} Verifica que $f(x) = 2x + 3$ y $g(x) = \dfrac{x - 3}{2}$ son funciones inversas.

  \item \medio{} Dadas $f(x) = \dfrac{1}{x + 1}$ y $g(x) = \dfrac{x}{x - 2}$, determina el dominio de $f \circ g$.

  \item \dificil{} Demuestra que $f(x) = x^2 - 4x$ (con dominio restringido $x \ge 2$) es inyectiva y halla su inversa.

  \item \dificil{} Halla la función inversa de $f(x) = \dfrac{2x - 1}{x + 3}$.

  \item \dificil{} Si $f(x) = x + 1$ y $(f \circ g)(x) = x^2 - 3$, determina $g(x)$.

  \item \muyDificil{} Halla la inversa de $f(x) = \dfrac{3x - 2}{2x + 1}$ y verifica que $(f \circ f^{-1})(x) = x$.
\end{enumerate}

\subsection{Ejercicios propuestos: Retos Integradores}

\begin{enumerate}[resume=ejercicios, leftmargin=*]
  \item \facil{} Determina si $f(x) = x^3$ es inyectiva y justifica tu respuesta.

  \item \medio{} Dadas $f(x) = x^2 + 1$ y $g(x) = \sqrt{x}$, halla $(f \circ g)(x)$ y su dominio.

  \item \medio{} Halla el rango de $f(x) = \dfrac{x + 2}{x - 3}$.

  \item \medio{} Dadas $f(x) = 2x$, $g(x) = x + 1$ y $h(x) = x^2$, halla $(f \circ g \circ h)(x)$.

  \item \dificil{} Sea $f(x) = \dfrac{x}{x + 1}$. Determina la expresión de $f\big(f(x)\big)$.

  \item \dificil{} Determina el número de asíntotas (verticales, horizontales u oblicuas) de $f(x) = \dfrac{2x^2 + 3x}{x^2 - 1}$.

  \item \muyDificil{} Halla todas las asíntotas de $f(x) = \dfrac{x^3 - 2x}{x^2 + 1}$.

  \item \muyDificil{} Determina $a$ y $b$ para que $f(x) = \dfrac{ax + b}{x - 1}$ cumpla $f(2) = 5$ y $f^{-1}(3) = 0$.
\end{enumerate}
\newpage
% ============================================================
\section{Clave de Respuestas}
% ============================================================

\subsection*{Sección 1: Funciones Racionales (Ejercicios 1--12)}

\begin{enumerate}[series=respuestas, leftmargin=*, label=\textbf{\arabic*.}]
  \item $(-\infty, 5) \cup (5, \infty)$.

  \item $(-\infty, -3) \cup (-3, 3) \cup (3, \infty)$.

  \item A.V.: $x = 2$; A.H.: $y = 0$.

  \item A.V.: $x = 0$; A.H.: $y = 5$.

  \item $f(x) = \dfrac{(x - 2)(x + 2)}{(x - 3)(x + 1)}$ (sin cancelaciones). Dominio: $x \neq 3$, $x \neq -1$. A.V.: $x = 3$ y $x = -1$. A.H.: $y = 1$ (mismos grados). Intercepto $y$: $(0, \frac{4}{-3}) = (0, -\frac{4}{3})$. Interceptos $x$: $x = \pm 2$: $(2, 0)$ y $(-2, 0)$.

  \item División: $2x^2 + 5x + 3$ entre $x - 1$ da $2x + 7 + \frac{10}{x - 1}$. A.O.: $y = 2x + 7$; A.V.: $x = 1$.

  \item No hay A.V. ($x^2 + 1 > 0$ siempre); A.H.: $y = 1$.

  \item $f(x) = \dfrac{(x - 2)(x^2 + 2x + 4)}{x - 2} = x^2 + 2x + 4$ para $x \neq 2$. Tiene un \textbf{hueco} en $(2, 12)$, no asíntota.

  \item Cociente: $3x + 2$ y residuo $-2x - 1$. Asíntota oblicua: $y = 3x + 2$.

  \item Grado numerador = grado denominador = 1. Asíntota horizontal: $y = \frac{k}{1} = k = 4$. Entonces $k = 4$.

  \item $f(x) = \dfrac{3(x + 1)}{x + 2}$ (verificar: A.V. en $x = -2$, A.H. en $y = 3$, intercepto $y$ en $\frac{3}{2}$).

  \item $f(x) = \dfrac{(x - 1)(x + 1)}{(x - 2)^2}$. A.V.: $x = 2$ (doble); A.H.: $y = 1$; intercepto $y$: $(0, -\frac{1}{4})$; interceptos $x$: $(-1, 0)$ y $(1, 0)$. Sin huecos.
\end{enumerate}

\newpage


\subsection*{Sección 2: Composición e Inversas (Ejercicios 13--24)}

\begin{enumerate}[resume=respuestas, leftmargin=*, label=\textbf{\arabic*.}]
  \item a) $4x + 2$ \quad b) $-2x + 2$ \quad c) $3x^2 + 6x$ \quad d) $\dfrac{x + 2}{3x}$ (con $x \neq 0$).

  \item $(f \circ g)(x) = (x + 1)^2 = x^2 + 2x + 1$; \quad $(g \circ f)(x) = x^2 + 1$.

  \item $f^{-1}(x) = \dfrac{x + 7}{3}$.

  \item $f^{-1}(x) = \sqrt[3]{x - 5}$.

  \item $(f \circ g)(x) = \sqrt{x^2 + 1}$. Dominio: $\mathbb{R}$ (pues $x^2 + 1 \geq 1 > 0$ para todo $x$).

  \item No es inyectiva en $\mathbb{R}$ (parábola). Se puede restringir a $[2, \infty)$ o $(-\infty, 2]$ para que sea inyectiva.

  \item $y = \frac{x}{x - 2}$. Intercambiamos: $x = \frac{y}{y - 2} \Rightarrow x(y - 2) = y \Rightarrow xy - 2x = y \Rightarrow xy - y = 2x \Rightarrow y(x - 1) = 2x \Rightarrow f^{-1}(x) = \dfrac{2x}{x - 1}$ (con $x \neq 1$).

  \item $f(g(x)) = \frac{3\left(\frac{2x + 1}{3}\right) - 1}{2} = \frac{2x + 1 - 1}{2} = x$. $g(f(x)) = \frac{2\left(\frac{3x - 1}{2}\right) + 1}{3} = \frac{3x - 1 + 1}{3} = x$. Sí son inversas.

  \item Completamos el cuadrado: $y = (x - 2)^2 - 1$. Intercambiamos: $x = (y - 2)^2 - 1 \Rightarrow (y - 2)^2 = x + 1$. Con dominio $y \geq 2$: $y - 2 \geq 0 \Rightarrow y = 2 + \sqrt{x + 1}$. $f^{-1}(x) = 2 + \sqrt{x + 1}$, dominio $x \geq -1$.

  \item Para $(f \circ g)(x) = f(g(x)) = \dfrac{1}{x^2 - 1}$: dominio $x \neq \pm 1$. Para $(g \circ f)(x) = g(f(x)) = \left(\dfrac{1}{x - 1}\right)^2$: dominio $x \neq 1$.

  \item $f(g(x)) = f(x + 3) = 2x - 1$. Sea $u = x + 3$, entonces $x = u - 3$: $f(u) = 2(u - 3) - 1 = 2u - 7$. Por tanto $f(x) = 2x - 7$.

  \item $y = 0{,}8x - 10$. Intercambiamos: $x = 0{,}8y - 10 \Rightarrow 0{,}8y = x + 10 \Rightarrow y = \dfrac{x + 10}{0{,}8} = 1{,}25x + 12{,}5$. $P^{-1}(x) = 1{,}25x + 12{,}5$ da el precio original antes de impuesto y descuento, dado el precio final.
\end{enumerate}


% ──────────────────────────────────────────────────────────
%  NUEVAS RESPUESTAS (26) — INSERCIÓN MANUAL
% ──────────────────────────────────────────────────────────
\subsection*{Sección 3: Asíntotas y Gráficas Racionales (Ejercicios 25--33)}
\begin{enumerate}[resume=respuestas, leftmargin=*, label=\textbf{\arabic*.}]
  \item Dominio: $\mathbb{R} \setminus \{\pm 3\}$
  \item Asíntota vertical: $x = 4$
  \item Asíntota vertical $x = 1$; \quad asíntota horizontal $y = 3$
  \item $y = \dfrac{2}{3}$
  \item Discontinuidad evitable (agujero) en $x = 2$; simplificada: $f(x) = x + 2$ (con $x \neq 2$).
  \item Asíntota vertical $x = 1$; \quad asíntota oblicua $y = x + 1$
  \item $y = 2x - 7$
  \item Dominio $\mathbb{R} \setminus \{\pm 2\}$; asíntotas verticales $x = \pm 2$; asíntota horizontal $y = 0$; intercepto en $(0, 0)$.
  \item Dominio $\mathbb{R} \setminus \{\pm 1\}$; asíntotas verticales $x = \pm 1$; asíntota oblicua $y = x$.
\end{enumerate}

\subsection*{Sección 4: Composición e Inversas Avanzadas (Ejercicios 34--42)}
\begin{enumerate}[resume=respuestas, leftmargin=*, label=\textbf{\arabic*.}]
  \item $(f \circ g)(x) = 2x^2 + 1$; \quad $(g \circ f)(x) = (2x + 1)^2 = 4x^2 + 4x + 1$
  \item $f^{-1}(x) = \dfrac{x + 6}{3}$
  \item $f^{-1}(x) = x^2 + 2$, con dominio $x \ge 0$ (rango de $f$)
  \item $(f \circ g)(x) = 2\left(\dfrac{x - 3}{2}\right) + 3 = x$ \quad y \quad $(g \circ f)(x) = \dfrac{(2x + 3) - 3}{2} = x$. Son inversas.
  \item $f(g(x)) = \dfrac{x - 2}{2(x - 1)}$; dominio: $\mathbb{R} \setminus \{1, 2\}$
  \item $y = x^2 - 4x \Rightarrow (x - 2)^2 = y + 4$; con $x \ge 2$: $f^{-1}(x) = 2 + \sqrt{x + 4}$, dominio $x \ge -4$.
  \item $f^{-1}(x) = \dfrac{3x + 1}{2 - x}$, dominio $\mathbb{R} \setminus \{2\}$
  \item $g(x) = x^2 - 4$ \quad (pues $f(g(x)) = g(x) + 1 = x^2 - 3$)
  \item $f^{-1}(x) = \dfrac{x + 2}{3 - 2x}$; se verifica $(f \circ f^{-1})(x) = x$ sustituyendo y simplificando.
\end{enumerate}

\subsection*{Sección 5: Retos Integradores (Ejercicios 43--50)}
\begin{enumerate}[resume=respuestas, leftmargin=*, label=\textbf{\arabic*.}]
  \item Sí es inyectiva: es estrictamente creciente en $\mathbb{R}$ (toda recta horizontal corta una sola vez).
  \item $(f \circ g)(x) = x + 1$, con dominio $x \ge 0$
  \item Rango: $\mathbb{R} \setminus \{1\}$ \quad (despejando $x = \dfrac{3y + 2}{y - 1}$)
  \item $(f \circ g \circ h)(x) = 2(x^2 + 1) = 2x^2 + 2$
  \item $f\big(f(x)\big) = \dfrac{x}{2x + 1}$ \quad (con $x \neq 0$ y $x \neq -1$)
  \item Dos asíntotas verticales ($x = \pm 1$) y una asíntota horizontal ($y = 2$); no hay oblicua.
  \item No hay asíntotas verticales ($x^2 + 1$ nunca se anula); asíntota oblicua $y = x$; no hay horizontal.
  \item $f(2) = 5 \Rightarrow 2a + b = 5$; \quad $f^{-1}(3) = 0 \Rightarrow f(0) = 3 \Rightarrow b = -3$; \quad luego $a = 4$. \quad $f(x) = \dfrac{4x - 3}{x - 1}$.
\end{enumerate}
\end{document}
